% ============================================================================
% APPENDIX GN: LIT-GNEEZY - Incentives, Gender & Deception
% ============================================================================
% Category: Literature Integration (LIT-)
% Language: English
% EBF Integration: Incentive design, gender differences, and honesty
% ============================================================================

\chapter*{Appendix UNMAPPED_GN: Uri Gneezy Research Integration}
\addcontentsline{toc}{chapter}{Appendix UNMAPPED_GN: LIT-GNEEZY - Incentives, Gender \& Deception}
\label{app:gneezy}

% ============================================================================
% HEADER BLOCK
% ============================================================================
\begin{tcolorbox}[colback=white,colframe=black!75!white,title=Appendix Metadata]
\textbf{Appendix:} GN -- LIT-GNEEZY \\
\textbf{Category:} Literature Integration (LIT-) \\
\textbf{Version:} 1.0 (January 2026) \\
\textbf{Dependencies:} CORE-HOW, LIT-FEHR, LIT-LIST, CONTEXT-DEMOGRAPHIC \\
\textbf{Status:} Complete
\end{tcolorbox}

% ============================================================================
% CROSS-REFERENCE MAP
% ============================================================================
\begin{tcolorbox}[colback=blue!5!white,colframe=blue!75!black,title=Cross-Reference Map]
\textbf{Builds On:}
\begin{itemize}
    \item Appendix K (LIT-FEHR): Social preferences and reciprocity
    \item Appendix AB (LIT-LIST): Field experiments methodology
    \item Appendix UNMAPPED_B (CORE-HOW): Complementarity and incentive interactions
    \item Appendix UNMAPPED_MM (LIT-MARECHAL): Honesty research
\end{itemize}

\textbf{Extends To:}
\begin{itemize}
    \item Appendix AG (DOMAIN-HR): Workplace incentive design
    \item Appendix AE (DOMAIN-HEALTH): Health behavior incentives
    \item Appendix AI (CONTEXT-DEMOGRAPHIC): Gender effects
    \item Appendix AO (PREDICT-BEHAVIOR): Behavioral prediction
\end{itemize}

\textbf{Methods Connection:}
\begin{itemize}
    \item Appendix E (METHOD-E): Field experimental methodology
    \item Appendix UNMAPPED_FI (LIT-FISCHBACHER): Laboratory methods
\end{itemize}
\end{tcolorbox}

% ============================================================================
% CHAPTER LINKAGE
% ============================================================================
\begin{tcolorbox}[colback=green!5!white,colframe=green!75!black,title=Chapter Linkage]
\textbf{Primary:} Chapter 9 (Incentives \& Motivation) \\
\textbf{Secondary:} Chapter 5 (Social Preferences), Chapter 14 (Organizational Applications)
\end{tcolorbox}

% ============================================================================
% ABSTRACT
% ============================================================================
\begin{abstract}
This appendix integrates Uri Gneezy's influential research on incentives, gender differences in competition, and deception into the Evidence-Based Framework. Gneezy's work fundamentally changed how economists understand the relationship between extrinsic incentives and intrinsic motivation, demonstrating that monetary incentives can backfire by ``crowding out'' prosocial behavior. His cross-cultural studies on gender and competition revealed that competitive behavior is substantially shaped by culture, not just biology. This appendix provides critical parameters for CORE-HOW (incentive complementarities) and establishes boundary conditions for when incentives help versus harm.
\end{abstract}

% ============================================================================
% QUICK REFERENCE
% ============================================================================
\begin{tcolorbox}[colback=gray!10!white,colframe=gray!50!black,title=Quick Reference]
Key terms defined in this appendix (see also Appendix G for complete Glossary):
\begin{description}
    \item[Crowding Out] When extrinsic incentives reduce intrinsic motivation
    \item[A Fine is a Price] Principle that penalties can legitimize otherwise stigmatized behavior
    \item[Competition Gap] Gender difference in willingness to enter competitive environments
    \item[Lying Aversion] Intrinsic disutility from telling lies, independent of consequences
\end{description}
\end{tcolorbox}

% ============================================================================
% FUNDAMENTAL QUESTION
% ============================================================================
% -----------------------------------------------------------------------------
% APPENDIX SCOPE BOX
% -----------------------------------------------------------------------------

\begin{tcolorbox}[
    colback=orange!5!white,
    colframe=orange!75!black,
    title={\textbf{Appendix Scope: Ziel / In-Scope / Out-of-Scope / Constraints}},
    fonttitle=\bfseries
]

\textbf{Ziel (Objective):}
\begin{itemize}[nosep]
    \item When do incentives enhance versus undermine desired behavior, and what role do social context and individual characteristics (like gender) play in moderating these effects?
\end{itemize}

\textbf{In-Scope:}
\begin{itemize}[nosep]
    \item Fundamental Question
    \item Theoretical Framework
    \item Core Axioms
    \item Key Empirical Results
    \item EBF Integration
\end{itemize}

\textbf{Out-of-Scope (delegiert an andere Appendices):}
\begin{itemize}[nosep]
    \item LIT-FEHR $\rightarrow$ Appendix K
    \item LIT-LIST $\rightarrow$ Appendix AB
    \item CORE-HOW $\rightarrow$ Appendix UNMAPPED_B
    \item LIT-MARECHAL $\rightarrow$ Appendix UNMAPPED_MM
    \item DOMAIN-HR $\rightarrow$ Appendix AG
\end{itemize}

\textbf{Constraints (Anwendungsgrenzen):}
\begin{itemize}[nosep]
    \item [TODO: Define when this appendix does NOT apply]
    \item [TODO: Define required prerequisites]
    \item [TODO: Define known limitations]
\end{itemize}

\textbf{Lieferobjekte (Deliverables):}
\begin{itemize}[nosep]
    \item 5 Table(s)
    \item 10 Equation(s)
    \item 6 Axiom(s)
    \item Worked Example(s)
\end{itemize}

\end{tcolorbox}


\section{Fundamental Question}

\begin{tcolorbox}[colback=yellow!10!white,colframe=orange!75!black,title=Central Research Question]
\textbf{When do incentives enhance versus undermine desired behavior, and what role do social context and individual characteristics (like gender) play in moderating these effects?}
\end{tcolorbox}

This question addresses CORE-HOW ($\gamma$): How do different motivational forces interact?

% ============================================================================
% THEORETICAL FRAMEWORK
% ============================================================================
\section{Theoretical Framework}

\subsection{The Incentive-Motivation Nexus}

Gneezy's research establishes that behavior results from the interaction of extrinsic and intrinsic motivation:

\begin{equation}
B = f(M_{extrinsic}, M_{intrinsic}, \gamma_{crowd})
\end{equation}

Where:
\begin{itemize}
    \item $B$ = Behavioral outcome
    \item $M_{extrinsic}$ = Extrinsic incentive (money, punishment)
    \item $M_{intrinsic}$ = Intrinsic motivation (social norms, identity)
    \item $\gamma_{crowd}$ = Crowding parameter ($<0$ = crowding out, $>0$ = crowding in)
\end{itemize}

\subsection{The ``Fine is a Price'' Mechanism}

When a fine is introduced for undesirable behavior:

\begin{equation}
U_{late} = \begin{cases}
-\text{Social Stigma} & \text{(before fine)} \\
-\text{Fine} + \text{Legitimacy Gain} & \text{(after fine)}
\end{cases}
\end{equation}

If $|\text{Legitimacy Gain}| > |\text{Fine}| - |\text{Social Stigma}|$, the fine \textit{increases} the undesirable behavior.

\subsection{Gender and Competition}

\begin{equation}
P(\text{compete})_i = \alpha + \beta_1 \cdot \text{Ability}_i + \beta_2 \cdot \text{Gender}_i + \beta_3 \cdot \text{Culture}_i + \epsilon_i
\end{equation}

Where:
\begin{itemize}
    \item $\beta_2 < 0$ for women in patriarchal societies
    \item $\beta_2 \approx 0$ in matrilineal societies (Khasi)
    \item Culture explains $>50\%$ of gender gap
\end{itemize}

% ============================================================================
% AXIOMS
% ============================================================================
\section{Core Axioms}

\begin{tcolorbox}[colback=red!5!white,colframe=red!75!black,title=Axiom UNMAPPED_GN-1: A Fine is a Price]
\textbf{Monetary penalties can increase undesirable behavior by transforming moral transgressions into market transactions.}

\begin{equation}
\frac{\partial B_{undesirable}}{\partial \text{Fine}} > 0 \quad \text{when} \quad \gamma_{crowd} < 0
\end{equation}

Daycare late pickups increased 100\% after fines were introduced.

\textit{Evidence:} Gneezy \& Rustichini (2000), $N=10$ daycares, Journal of Legal Studies.
\end{tcolorbox}

\begin{tcolorbox}[colback=red!5!white,colframe=red!75!black,title=Axiom UNMAPPED_GN-2: Gender Competition Gap]
\textbf{Women are less likely to self-select into competitive environments, even when equally capable.}

\begin{equation}
P(\text{tournament} | \text{Female}) \approx 0.35 \quad \text{vs.} \quad P(\text{tournament} | \text{Male}) \approx 0.73
\end{equation}

The 38 percentage point gap persists controlling for ability and risk preferences.

\textit{Evidence:} Niederle \& Gneezy (2007), QJE.
\end{tcolorbox}

\begin{tcolorbox}[colback=red!5!white,colframe=red!75!black,title=Axiom UNMAPPED_GN-3: Cultural Malleability of Gender Gap]
\textbf{The gender competition gap is substantially determined by culture, not biology.}

\begin{equation}
\text{Gap}_{Khasi} \approx -0.08 \quad \text{vs.} \quad \text{Gap}_{Masai} \approx +0.50
\end{equation}

In matrilineal Khasi society, women compete \textit{more} than men.

\textit{Evidence:} Gneezy, Leonard \& List (2009), Econometrica.
\end{tcolorbox}

\begin{tcolorbox}[colback=red!5!white,colframe=red!75!black,title=Axiom UNMAPPED_GN-4: Lying Aversion]
\textbf{People have intrinsic costs of lying that are sensitive to both the size and direction of the lie.}

\begin{equation}
U_{lie} = G_{material} - c(\text{lie size}) - \phi \cdot \mathbf{1}_{harm\_other}
\end{equation}

Where $c(\cdot)$ is convex (larger lies are disproportionately costly) and $\phi > 0$.

\textit{Evidence:} Gneezy (2005), AER; Gneezy, Kajackaite \& Sobel (2017), AER.
\end{tcolorbox}

\begin{tcolorbox}[colback=red!5!white,colframe=red!75!black,title=Axiom UNMAPPED_GN-5: High Stakes Choking]
\textbf{Very high incentives can impair performance on cognitively demanding tasks.}

\begin{equation}
\text{Performance} = \alpha + \beta \cdot \text{Incentive} - \gamma \cdot \text{Incentive}^2
\end{equation}

The inverted-U relationship peaks at moderate incentive levels.

\textit{Evidence:} Ariely, Gneezy, Loewenstein \& Mazar (2009), REStud.
\end{tcolorbox}

\begin{tcolorbox}[colback=red!5!white,colframe=red!75!black,title=Axiom UNMAPPED_GN-6: Lab-Field Gap]
\textbf{Laboratory measures of social preferences may overstate prosociality relative to field behavior.}

\begin{equation}
\text{Prosociality}_{lab} > \text{Prosociality}_{field}
\end{equation}

Scrutiny, selection, and context drive the gap.

\textit{Evidence:} Levitt, List \& Gneezy (2007), JEP.
\end{tcolorbox}

% ============================================================================
% KEY EMPIRICAL RESULTS
% ============================================================================
\section{Key Empirical Results}

\subsection{The Daycare Study (2000)}

\begin{table}[h]
\centering
\begin{tabular}{lcc}
\hline
\textbf{Period} & \textbf{Late Pickups/Week} & \textbf{Change} \\
\hline
Before fine & 8 & -- \\
After fine introduced & 16 & +100\% \\
After fine removed & 16 & No recovery \\
\hline
\end{tabular}
\caption{Late pickups in Israeli daycares}
\end{table}

\textbf{Key insight:} Once social norms are replaced by market norms, they don't return.

\subsection{Gender Competition Studies}

\begin{table}[h]
\centering
\begin{tabular}{lccc}
\hline
\textbf{Society} & \textbf{Men Compete} & \textbf{Women Compete} & \textbf{Gap} \\
\hline
US (lab) & 73\% & 35\% & +38pp \\
Masai (patriarchal) & 50\% & 26\% & +24pp \\
Khasi (matrilineal) & 39\% & 54\% & -15pp \\
\hline
\end{tabular}
\caption{Competition choices across cultures}
\end{table}

\textbf{Key insight:} Biology alone cannot explain the gender gap.

\subsection{Deception and Consequences}

\begin{table}[h]
\centering
\begin{tabular}{lcc}
\hline
\textbf{Condition} & \textbf{Lies (\%)} & \textbf{Interpretation} \\
\hline
Lie helps self, hurts other & 36\% & Baseline deception \\
Lie helps self, helps other & 55\% & Less aversion when both gain \\
Lie helps self only & 17\% & Aversion to pure selfishness \\
\hline
\end{tabular}
\caption{Lying rates by consequence structure}
\end{table}

% ============================================================================
% EBF INTEGRATION
% ============================================================================
\section{EBF Integration}

\subsection{CORE-HOW Specification}

Gneezy's research provides critical parameters for the complementarity function:

\begin{equation}
\gamma_{incentive} = \begin{cases}
> 0 & \text{(market context, clear exchange)} \\
< 0 & \text{(social context, moral framing)} \\
\approx 0 & \text{(neutral context)}
\end{cases}
\end{equation}

\begin{table}[h]
\centering
\begin{tabular}{lccl}
\hline
\textbf{Parameter} & \textbf{Estimate} & \textbf{SE} & \textbf{Source} \\
\hline
$\gamma_{crowd}$ (social context) & -0.50 & 0.15 & Gneezy \& Rustichini (2000) \\
Gender gap (US) & 0.38 & 0.05 & Niederle \& Gneezy (2007) \\
Gender gap (matrilineal) & -0.15 & 0.08 & Gneezy et al. (2009) \\
Lying cost per \$1 & 0.12 & 0.03 & Gneezy (2005) \\
\hline
\end{tabular}
\caption{EBF parameters from Gneezy research}
\end{table}

\subsection{10C Framework Mapping}

\begin{table}[h]
\centering
\begin{tabular}{lll}
\hline
\textbf{CORE} & \textbf{Gneezy Contribution} & \textbf{Key Paper} \\
\hline
WHO (AAA) & Self vs. others in deception & Gneezy (2005) \\
WHAT (C) & Intrinsic vs. extrinsic utility & Gneezy \& Rustichini (2000) \\
HOW (B) & Crowding out complementarity & All incentive work \\
WHEN (V) & Market vs. social framing & Fine is a Price \\
WHERE (BBB) & $\gamma_{crowd}$ parameters & Meta-analysis \\
AWARE (AU) & Lab-field awareness gap & Levitt et al. (2007) \\
\hline
\end{tabular}
\caption{10C Framework mapping for Gneezy research}
\end{table}

% ============================================================================
% WORKED EXAMPLE
% ============================================================================
\section{Worked Example: Designing Employee Wellness Incentives}

\begin{tcolorbox}[colback=blue!5!white,colframe=blue!75!black,title=Application: Avoiding Crowding Out in Health Programs]
\textbf{Context:} A company wants to increase gym attendance through incentives.

\textbf{Step 1: Diagnose Current Motivation}
\begin{itemize}
    \item Current attendance: 15\% of employees
    \item Primary motivation: Health identity, social connection
    \item Risk: Monetary incentives may crowd out intrinsic motivation
\end{itemize}

\textbf{Step 2: Apply Gneezy Framework}

Using crowding parameters:
\begin{itemize}
    \item If framed as payment: $\gamma_{crowd} < 0$ (backfire risk)
    \item If framed as recognition: $\gamma_{crowd} > 0$ (enhancement)
\end{itemize}

\textbf{Step 3: Design Intervention}

\begin{enumerate}
    \item \textbf{Avoid direct cash:} Use non-monetary recognition
    \item \textbf{Frame as support:} ``Company invests in your health''
    \item \textbf{Social component:} Team challenges, not individual bonuses
    \item \textbf{Commitment device:} Pre-commitment with social accountability
\end{enumerate}

\textbf{Step 4: Gender Considerations}

Per Gneezy gender research:
\begin{itemize}
    \item Avoid competitive framing (tournament structures)
    \item Offer choice between competitive and collaborative
    \item Expected: 20\% higher female participation vs. tournament-only
\end{itemize}

\textbf{Step 5: Predict Effects}

\begin{align}
\text{Attendance (cash bonus)} &\approx 15\% + 5\% - 8\%_{crowd} = 12\% \text{ (worse)} \\
\text{Attendance (team challenge)} &\approx 15\% + 10\% + 3\%_{enhance} = 28\% \text{ (better)}
\end{align}
\end{tcolorbox}

% ============================================================================
% IMPLICATIONS
% ============================================================================
\section{Implications for Practice}

\subsection{For Incentive Design}

\begin{enumerate}
    \item \textbf{Context First:} Diagnose whether intrinsic motivation exists before adding incentives
    \item \textbf{Frame Carefully:} Gift framing beats payment framing
    \item \textbf{Avoid Small Fines:} Either large penalties or social norms---not weak fines
    \item \textbf{Consider Removal:} Once norms are displaced, they may not return
\end{enumerate}

\subsection{For Gender Equity}

\begin{enumerate}
    \item \textbf{Offer Choice:} Let employees choose competitive or non-competitive tracks
    \item \textbf{Reduce Stakes:} Lower-stakes competition reduces gap
    \item \textbf{Team Settings:} Mixed-gender teams reduce individual competition aversion
    \item \textbf{Culture Change:} Long-term norm shifts are possible (Khasi evidence)
\end{enumerate}

% ============================================================================
% SUMMARY
% ============================================================================

% === AUTO-GENERATED ENTRIES (2026-02-22) ===

%%%%%%%%%%%%%%%%%%%%%%%%%%%%%%%%%%%%%%%%%%%%%%%%%%%%%%%%%%%%%%%%%%%%%%%%%%%%%%%
\subsubsection{2000: Pay Enough or Don't Pay at All}
%%%%%%%%%%%%%%%%%%%%%%%%%%%%%%%%%%%%%%%%%%%%%%%%%%%%%%%%%%%%%%%%%%%%%%%%%%%%%%%

\paragraph{Citation.}
Gneezy, Uri and Rustichini, Aldo (2000). ``Pay Enough or Don't Pay at All.'' \textit{Quarterly Journal of Economics}, 115, 791--810.

\paragraph{10C Integration.}
\begin{itemize}[nosep]
  \item \textbf{Domain}: [To be specified]
  \item \textbf{use\_for}: LIT-GNEEZY, LIT-O
\end{itemize}


%%%%%%%%%%%%%%%%%%%%%%%%%%%%%%%%%%%%%%%%%%%%%%%%%%%%%%%%%%%%%%%%%%%%%%%%%%%%%%%
\subsubsection{2003: Bargaining under a Deadline: Evidence from the Reverse Ultim...}
%%%%%%%%%%%%%%%%%%%%%%%%%%%%%%%%%%%%%%%%%%%%%%%%%%%%%%%%%%%%%%%%%%%%%%%%%%%%%%%

\paragraph{Citation.}
Gneezy, Uri and Haruvy, Ernan and Roth, Al (2003). ``Bargaining under a Deadline: Evidence from the Reverse Ultimatum Game.'' \textit{Games and Economic Behavior}, 45.

\paragraph{10C Integration.}
\begin{itemize}[nosep]
  \item \textbf{Domain}: [To be specified]
  \item \textbf{use\_for}: LIT-GNEEZY, LIT-O
\end{itemize}


%%%%%%%%%%%%%%%%%%%%%%%%%%%%%%%%%%%%%%%%%%%%%%%%%%%%%%%%%%%%%%%%%%%%%%%%%%%%%%%
\subsubsection{2003: Performance in Competitive Environments: Gender Differences}
%%%%%%%%%%%%%%%%%%%%%%%%%%%%%%%%%%%%%%%%%%%%%%%%%%%%%%%%%%%%%%%%%%%%%%%%%%%%%%%

\paragraph{Citation.}
Gneezy, Uri and Niederle, Muriel and Rustichini, Aldo (2003). ``Performance in Competitive Environments: Gender Differences.'' \textit{null}, null.

\paragraph{10C Integration.}
\begin{itemize}[nosep]
  \item \textbf{Domain}: [To be specified]
  \item \textbf{use\_for}: LIT-GNEEZY, LIT-O
\end{itemize}


%%%%%%%%%%%%%%%%%%%%%%%%%%%%%%%%%%%%%%%%%%%%%%%%%%%%%%%%%%%%%%%%%%%%%%%%%%%%%%%
\subsubsection{2004: Incentives, Punishment, and Behavior}
%%%%%%%%%%%%%%%%%%%%%%%%%%%%%%%%%%%%%%%%%%%%%%%%%%%%%%%%%%%%%%%%%%%%%%%%%%%%%%%

\paragraph{Citation.}
Gneezy, Uri and Rustichini, Aldo (2004). ``Incentives, Punishment, and Behavior.'' \textit{Handbook of Experimental Economics}, 2.

\paragraph{10C Integration.}
\begin{itemize}[nosep]
  \item \textbf{Domain}: [To be specified]
  \item \textbf{use\_for}: LIT-GNEEZY, LIT-O
\end{itemize}


%%%%%%%%%%%%%%%%%%%%%%%%%%%%%%%%%%%%%%%%%%%%%%%%%%%%%%%%%%%%%%%%%%%%%%%%%%%%%%%
\subsubsection{2006: The Uncertainty Effect: When a Risky Prospect Is Valued Less...}
%%%%%%%%%%%%%%%%%%%%%%%%%%%%%%%%%%%%%%%%%%%%%%%%%%%%%%%%%%%%%%%%%%%%%%%%%%%%%%%

\paragraph{Citation.}
Gneezy, Uri and List, John and Wu, George (2006). ``The Uncertainty Effect: When a Risky Prospect Is Valued Less Than Its Worst Outcome.'' \textit{Quarterly Journal of Economics}, 121.

\paragraph{10C Integration.}
\begin{itemize}[nosep]
  \item \textbf{Domain}: [To be specified]
  \item \textbf{use\_for}: LIT-GNEEZY, LIT-O
\end{itemize}


%%%%%%%%%%%%%%%%%%%%%%%%%%%%%%%%%%%%%%%%%%%%%%%%%%%%%%%%%%%%%%%%%%%%%%%%%%%%%%%
\subsubsection{2012: Shared Social Responsibility: A Field Experiment in Pay-What...}
%%%%%%%%%%%%%%%%%%%%%%%%%%%%%%%%%%%%%%%%%%%%%%%%%%%%%%%%%%%%%%%%%%%%%%%%%%%%%%%

\paragraph{Citation.}
Gneezy, Ayelet and Gneezy, Uri and Nelson, Leif and Brown, Amber (2012). ``Shared Social Responsibility: A Field Experiment in Pay-What-You-Want Pricing and Charitable Giving.'' \textit{Science}, 329.

\paragraph{10C Integration.}
\begin{itemize}[nosep]
  \item \textbf{Domain}: [To be specified]
  \item \textbf{use\_for}: LIT-GNEEZY, LIT-O
\end{itemize}


%%%%%%%%%%%%%%%%%%%%%%%%%%%%%%%%%%%%%%%%%%%%%%%%%%%%%%%%%%%%%%%%%%%%%%%%%%%%%%%
\subsubsection{2019: Bribery: Greed versus Reciprocity}
%%%%%%%%%%%%%%%%%%%%%%%%%%%%%%%%%%%%%%%%%%%%%%%%%%%%%%%%%%%%%%%%%%%%%%%%%%%%%%%

\paragraph{Citation.}
Gneezy, Uri and Saccardo, Silvia and Van Veldhuizen, Roel (2019). ``Bribery: Greed versus Reciprocity.'' \textit{Management Science}, 65.

\paragraph{10C Integration.}
\begin{itemize}[nosep]
  \item \textbf{Domain}: [To be specified]
  \item \textbf{use\_for}: LIT-GNEEZY, LIT-O
\end{itemize}


%%%%%%%%%%%%%%%%%%%%%%%%%%%%%%%%%%%%%%%%%%%%%%%%%%%%%%%%%%%%%%%%%%%%%%%%%%%%%%%
\subsubsection{2020: The Why Axis: Hidden Motives and the Undiscovered Economics ...}
%%%%%%%%%%%%%%%%%%%%%%%%%%%%%%%%%%%%%%%%%%%%%%%%%%%%%%%%%%%%%%%%%%%%%%%%%%%%%%%

\paragraph{Citation.}
Gneezy, Uri and List, John (2020). \textit{The Why Axis: Hidden Motives and the Undiscovered Economics of Everyday Life}. Publisher.

\paragraph{10C Integration.}
\begin{itemize}[nosep]
  \item \textbf{Domain}: [To be specified]
  \item \textbf{use\_for}: LIT-GNEEZY, LIT-O
\end{itemize}


%%%%%%%%%%%%%%%%%%%%%%%%%%%%%%%%%%%%%%%%%%%%%%%%%%%%%%%%%%%%%%%%%%%%%%%%%%%%%%%
\subsubsection{1972: On the Concept of Health Capital}
%%%%%%%%%%%%%%%%%%%%%%%%%%%%%%%%%%%%%%%%%%%%%%%%%%%%%%%%%%%%%%%%%%%%%%%%%%%%%%%

\paragraph{Citation.}
Grossman, M. (1972). ``On the Concept of Health Capital.'' \textit{Journal of Political Economy}.

\paragraph{Core Finding.}
Auto-imported from extracted_papers.yaml (GB_LIT-BECKER_human_capital.tex)

\paragraph{10C Integration.}
\begin{itemize}[nosep]
  \item \textbf{Domain}: [To be specified]
  \item \textbf{use\_for}: LIT-GNEEZY
\end{itemize}

% === END AUTO-GENERATED ENTRIES ===
\section{Summary}

\begin{tcolorbox}[colback=gray!10!white,colframe=gray!50!black,title=Key Takeaways]
\begin{enumerate}
    \item \textbf{A Fine is a Price:} Penalties can legitimize bad behavior (+100\% late pickups)
    \item \textbf{Crowding Out:} Extrinsic incentives can destroy intrinsic motivation
    \item \textbf{Gender Gap:} 38pp competition gap is largely cultural, not biological
    \item \textbf{Lying Aversion:} People have intrinsic costs of deception
    \item \textbf{High Stakes Hurt:} Very large incentives impair cognitive performance
    \item \textbf{Lab $\neq$ Field:} Social preferences differ across contexts
\end{enumerate}

\textbf{Core Insight:} Incentives interact with existing motivations---design for the whole system.
\end{tcolorbox}

% ============================================================================
% GLOSSARY
% ============================================================================
\section{Glossary}

Key terms defined in this appendix (see also Appendix G for complete Glossary):

\begin{description}
    \item[Crowding Out] Reduction in intrinsic motivation due to extrinsic incentives
    \item[Crowding In] Enhancement of intrinsic motivation due to appropriate incentives
    \item[Market Norms] Transactional framing where payment signals exchange
    \item[Social Norms] Relational framing where behavior signals identity
    \item[Competition Gap] Gender difference in willingness to compete
    \item[Lying Aversion] Intrinsic disutility from dishonesty
    \item[Pay What You Want] Pricing mechanism leaving payment to buyer discretion
\end{description}

% ============================================================================
% CRITICAL FOUNDATIONS
% ============================================================================
\section{Critical Foundations}

\begin{enumerate}
    \item \textbf{Deci (1971):} Cognitive evaluation theory and intrinsic motivation
    \item \textbf{Frey \& Oberholzer-Gee (1997):} Motivation crowding theory
    \item \textbf{Bénabou \& Tirole (2003):} Intrinsic and extrinsic motivation model
    \item \textbf{Fehr \& Falk (2002):} Psychological foundations of incentives
    \item \textbf{Bowles \& Polanía-Reyes (2012):} Economic incentives and social preferences
\end{enumerate}

% ============================================================================
% OPEN ISSUES
% ============================================================================
\section{Open Issues}

\begin{enumerate}
    \item \textbf{Crowding Prediction:} When exactly will crowding out occur?
    \item \textbf{Norm Recovery:} Can displaced social norms be restored?
    \item \textbf{Gender Intervention:} What interventions reduce the competition gap?
    \item \textbf{Optimal Incentive Size:} Where is the performance-incentive peak?
    \item \textbf{Cross-Cultural Stability:} How stable are Khasi/Masai findings over time?
\end{enumerate}

% ============================================================================
% REFERENCES
% ============================================================================
\section{References}

See Master Bibliography (bcm2\_master\_references.bib) for complete citations.

\nocite{bcm_master}

\textbf{Primary Gneezy Sources} (17 papers integrated):

\begin{itemize}
    \item Gneezy, U. \& Rustichini, A. (2000). A Fine is a Price. \textit{Journal of Legal Studies}, 29(1).
    \item Gneezy, U. (2005). Deception: The Role of Consequences. \textit{AER}, 95(1).
    \item Niederle, M. \& Gneezy, U. (2007). Do Women Shy Away from Competition? \textit{QJE}, 122(3).
    \item Gneezy, U., Leonard, K. \& List, J. (2009). Gender Differences in Competition. \textit{Econometrica}, 77(5).
    \item Gneezy, U., Meier, S. \& Rey-Biel, P. (2011). When and Why Incentives Don't Work. \textit{JEP}, 25(4).
    \item Ariely, D., Gneezy, U., Loewenstein, G. \& Mazar, N. (2009). Large Stakes and Big Mistakes. \textit{REStud}, 76(2).
    \item Gneezy, U., Kajackaite, A. \& Sobel, J. (2017). Lying Aversion and the Size of the Lie. \textit{AER}, 108(2).
    \item Levitt, S., List, J. \& Gneezy, U. (2007). What Do Lab Experiments Reveal? \textit{JEP}, 21(2).
\end{itemize}

\end{document}
